% !TeX program = pdflatex
\documentclass[11pt,a4paper]{article}
\usepackage{../../recursos/latex/estilo-notas}

\begin{document}
\cabecalhoaula{05}{Árvores de Regressão e \emph{Ensembles}}{AME §4.8--4.10; ISLP cap.~8}

%==============================================================================
\section*{Objetivos da aula}
%==============================================================================
Ao final desta aula o aluno deve ser capaz de:
\begin{itemize}[leftmargin=1.4cm]
  \item descrever como uma \textbf{árvore de regressão} particiona o espaço de
        covariáveis e prediz por médias locais;
  \item explicar a construção em duas etapas: \emph{crescimento} por divisões
        binárias recursivas e \emph{poda} por custo--complexidade;
  \item justificar, via redução de variância, por que \textbf{agregar} estimadores
        ajuda;
  \item descrever \textbf{bagging} e explicar como as \textbf{florestas
        aleatórias} \emph{descorrelacionam} as árvores, e o papel de $m$.
\end{itemize}

\begin{emsala}
Árvores são o método mais \emph{interpretável} que veremos --- e, sozinhas, das
piores em predição. A história da aula é como transformá-las, via agregação, em
um dos métodos \emph{mais poderosos} que existem (as florestas). Pergunta de
abertura: \emph{``como combinar muitos modelos ruins para obter um ótimo?''}
\end{emsala}

%==============================================================================
\section{Árvores de regressão}
%==============================================================================
Uma árvore parte o espaço de covariáveis em regiões disjuntas $R_1,\dots,R_J$ por
\emph{particionamento recursivo} e prediz, em cada região, a média das respostas de
treino ali contidas:
\begin{equation}\label{eq:arvore}
  g(\x) = \frac{1}{\abs{\{i:\X_i\in R_k\}}}\sum_{i:\X_i\in R_k} Y_i,
  \qquad \text{para } \x\in R_k .
\end{equation}
Para prever, percorre-se a árvore do topo às folhas seguindo as condições
(``$x_j<t$?'') em cada nó.

\begin{ideia}
Vantagens estruturais das árvores: (i) altamente \textbf{interpretáveis};
(ii) lidam \emph{nativamente} com covariáveis discretas (cortes do tipo
$x_j\in\{\text{verde},\text{vermelho}\}$ --- sem precisar criar \emph{dummies});
(iii) fazem \textbf{seleção de variáveis} automática; (iv) capturam
\textbf{interações} sem que precisemos especificá-las.
\end{ideia}

\subsection{Etapa 1 --- crescimento (divisões binárias recursivas)}
Idealmente escolheríamos a partição que minimiza o EQM
$P(T)=\sum_R\sum_{i:\X_i\in R}(Y_i-\widehat y_R)^2/n$. Mas buscar a árvore ótima é
computacionalmente inviável (NP-difícil). Usa-se então uma heurística
\emph{gananciosa}: a cada passo, escolhe-se a covariável $x_j$ e o corte $t$ que
mais reduzem a soma de quadrados ao dividir uma região em
$R_1=\{x_j<t\}$ e $R_2=\{x_j\ge t\}$:
\begin{equation}\label{eq:split}
  \min_{j,\,t}\ \sum_{i:\X_i\in R_1}(Y_i-\widehat y_{R_1})^2
            + \sum_{i:\X_i\in R_2}(Y_i-\widehat y_{R_2})^2 .
\end{equation}
Repete-se recursivamente em cada região-filha, até um critério de parada (ex.: menos
de 5 observações por folha). Isso gera uma árvore grande $T_0$, que ajusta bem o
treino mas tende a \textbf{superajustar}.

\begin{emsala}
Pergunta: \emph{``a árvore gananciosa acha a solução de
$y=\operatorname{sinal}(x_1x_2x_3)$?''} Oito regiões resolvem exatamente, e caberiam
em profundidade~3 --- então a resposta intuitiva é sim. A §3 da \texttt{Aula prática
05} mede: o melhor corte da raiz reduz $0{,}194\%$ do RSS dali, e nos dois filhos que
ele produz os melhores cortes reduzem $0{,}447\%$ e $0{,}942\%$; e o RSS
de teste fica \textbf{acima} do de chutar a média até a profundidade~5 --- pior a cada
nível --- caindo a $13\%$ dele só com a árvore cheia, de 59 folhas. Vale desenhar os
octantes no quadro e mostrar que qualquer corte deixa os dois lados com média zero.

E o corolário que interessa para a próxima subseção: com
\texttt{min\_impurity\_decrease=0,005} a árvore fica com \textbf{uma folha}. A parada
precoce não é uma poda mais barata --- é uma poda que decide antes de ter informação.
\end{emsala}

\subsection{Etapa 2 --- poda por custo--complexidade}
Para reduzir a variância, \emph{podamos} $T_0$. A poda por custo--complexidade
introduz um parâmetro $\alpha\ge 0$ e busca a subárvore $T\subseteq T_0$ que
minimiza
\begin{equation}\label{eq:poda}
  \sum_{k=1}^{\abs{T}}\sum_{i:\X_i\in R_k}(Y_i-\widehat y_{R_k})^2
   + \alpha\,\abs{T},
\end{equation}
onde $\abs{T}$ é o número de folhas. O termo $\alpha\abs{T}$ penaliza árvores
grandes --- é o mesmo espírito da regularização da Aula~02 (trocamos viés por menos
variância). $\alpha$ é o \textbf{tuning parameter}, escolhido por
\textbf{validação cruzada} (Aula~03): $\alpha=0$ devolve $T_0$;
$\alpha\to\infty$ colapsa para a raiz (média global).

\begin{atencao}
Crescer e podar usam objetivos diferentes \emph{de propósito}: cresce-se de forma
gananciosa (olhando só um passo à frente) e corrige-se globalmente na poda. Não
confunda profundidade máxima (parâmetro de crescimento) com $\alpha$ (parâmetro de
poda) --- ambos controlam complexidade, mas em etapas distintas.
\end{atencao}

%==============================================================================
\section{Por que agregar? Redução de variância}
%==============================================================================
Árvores isoladas têm \emph{alta variância}: pequenas mudanças nos dados mudam
bastante a árvore. A ideia dos \emph{ensembles} é \textbf{combinar} muitas funções
para reduzir essa variância.

\begin{proposicao}[Média de preditores reduz o risco]\label{prop:media}
Sejam $g_1,\dots,g_B$ estimadores \emph{não-viesados}
($\E[g_b(\x)]=r(\x)$), de mesma variância $v(\x)=\Var(g_b(\x))$ e
\emph{não-correlacionados}. Então a média $\bar g=\frac1B\sum_b g_b$ tem risco
\[
  \E[(Y-\bar g(\x))^2\mid\x] = \Var(Y\mid\x) + \frac{v(\x)}{B}
  \;\le\; \E[(Y-g_b(\x))^2\mid\x].
\]
\end{proposicao}
\begin{proof}
Pela decomposição viés--variância (Aula~01), o risco de qualquer preditor $g$ é
$\Var(Y\mid\x)+\mathrm{viés}^2+\Var(g(\x))$. Como cada $g_b$ é não-viesado, $\bar
g$ também é (viés $=0$). Para a variância, sendo os $g_b$ não-correlacionados e de
variância $v(\x)$,
$\Var(\bar g(\x))=\frac1{B^2}\sum_b \Var(g_b(\x))=v(\x)/B$. Substituindo obtém-se a
igualdade; como $v(\x)/B\le v(\x)$, o risco não aumenta.
\end{proof}

\begin{ideia}
Duas exigências da Prop.~\ref{prop:media} guiam tudo o que segue: precisamos de
preditores (i) \textbf{não-viesados} --- por isso usamos árvores \emph{não podadas}
(profundas, baixo viés) --- e (ii) \textbf{pouco correlacionados} --- o ponto fraco,
que as florestas aleatórias atacam diretamente. Quanto menor a correlação, mais a
variância cai.
\end{ideia}

%==============================================================================
\section{Bagging}
%==============================================================================
\emph{Bagging} (\emph{bootstrap aggregating}) gera diversidade reamostrando os
dados. Para $b=1,\dots,B$: sorteie uma amostra \textbf{bootstrap} (com reposição,
tamanho $n$) e ajuste nela uma árvore \emph{profunda} (não podada) $g_b$. A predição
é a média:
\begin{equation}\label{eq:bagging}
  g(\x) = \frac1B\sum_{b=1}^B g_b(\x).
\end{equation}

%==============================================================================
\section{Florestas aleatórias}
%==============================================================================
O problema do \emph{bagging}: as árvores ficam \textbf{correlacionadas}. Se há uma
covariável muito forte, quase todas as árvores a usam no topo e ficam parecidas ---
e a Prop.~\ref{prop:media} rende pouco quando $\mathrm{Cor}(g_b)$ é alta.

\begin{ideia}
\textbf{Florestas aleatórias} (Breiman) \emph{descorrelacionam} as árvores: em
\emph{cada nó}, sorteia-se um subconjunto de apenas $m<p$ covariáveis e a divisão
só pode usar uma delas. Isso impede que a variável dominante apareça sempre,
diversificando as árvores e reduzindo mais a variância.
\end{ideia}

\noindent Detalhes:
\begin{itemize}[leftmargin=1.4cm,nosep]
  \item $m$ é o \emph{tuning parameter}; regra empírica $m\approx p/3$ em regressão
        (e $m\approx\sqrt p$ em classificação). $m=p$ recupera o \emph{bagging}.
  \item O risco é \textbf{robusto a $B$}: ele estabiliza e \emph{não há
        superajuste} ao aumentar $B$ --- não vale a pena fazer CV em $B$, basta
        usar $B$ ``grande o bastante''.
  \item A taxa de convergência das florestas depende
        essencialmente do número de \emph{variáveis relevantes} --- elas se adaptam
        à \emph{esparsidade}.
\end{itemize}

%==============================================================================
\section*{Resumo}
%==============================================================================
\begin{itemize}[leftmargin=1.4cm]
  \item \textbf{Árvore}: partição recursiva \eqref{eq:arvore}; cresce
        gananciosamente \eqref{eq:split} e poda por custo--complexidade
        \eqref{eq:poda} ($\alpha$ por CV). Interpretável, mas alta variância.
  \item Agregar reduz variância se os preditores são não-viesados e pouco
        correlacionados (Prop.~\ref{prop:media}).
  \item \textbf{Bagging} \eqref{eq:bagging}: árvores profundas em amostras
        bootstrap, agregadas pela média.
  \item \textbf{Florestas}: \emph{descorrelacionam} sorteando $m\approx p/3$
        variáveis por nó; robustas a $B$; adaptam-se à esparsidade.
\end{itemize}

\paragraph{Leitura recomendada.} [AME] §4.8 (árvores) e §4.9 (\emph{bagging} e
florestas); [ISLP] §8.1 e §8.2 (\emph{Tree-Based Methods}).

\end{document}
